\pdfoutput=1

\documentclass{article}

\usepackage[preprint]{tmlr}

\usepackage{amsmath,amssymb,amsfonts}
\usepackage{booktabs}
\usepackage{graphicx}
\usepackage{hyperref}
\usepackage{xcolor}
\usepackage{multirow}
\usepackage{microtype}

\renewcommand{\arraystretch}{1.3}
\emergencystretch=1.5em

\title{Validating the Lookback Mechanism: Framing Dissociation Confirms Belief-Specific Subspaces}

\author{\name Elliot Tower \\
      \email elliot@elliottower.ai}

\begin{document}
\maketitle

\begin{abstract}
Causal findings in mechanistic interpretability are described with specific labels (belief tracking, planning, deception) whose specificity outruns what the causal evidence can verify. \citet{prakash2025lookback} report a ``lookback mechanism'' in Llama-3.1-70B-Instruct that tracks characters' beliefs through pointer dereference via attention, identifying low-rank subspaces at three layer ranges that carry belief content when swapped between stories. We apply 35 criteria from the Mechanistic Validity framework and identify six untested assumptions, most critically that the evaluation never tests false beliefs: the story-generation code sets each character's belief equal to reality in every pair. We pre-register 15 experiments and complete nine. The subspaces are real: they beat shuffled-label and random-subspace baselines (all $p < 0.05$). They generalize to genuine false-belief stories (IIA $= 0.975$) and observability-mediated knowledge (IIA $= 0.988$), overturning our pre-registered prediction that they would not. The decisive result is a framing dissociation: the same story pairs produce full subspace transfer under belief questions (IIA $= 1.0$), partial transfer under action-recall questions, and zero transfer when the question asks about reality or uses fill-completion framing---even though the underlying fact is identical. The subspaces encode character-perspective information, not generic entity-state lookup. E-value analysis finds 19 of 20 positive effects robust to unmeasured confounding ($E \geq 2.0$). A cross-model mismatch test on Qwen2.5-14B-Instruct confirms the answer lookback's dependency structure transfers (singles near zero, pair at $1.0$) while the binding mechanism shows a structurally different pattern---partial transport that full-residual patching cannot detect. We upgrade the claim from ``Causally Suggestive'' to ``Mechanistically Supported'' and propose the term \emph{belief-constitutive entity binding}.
\end{abstract}

\noindent\textbf{Keywords:} mechanistic interpretability, validity, belief tracking, interchange interventions, replication


\section{Introduction}
\label{sec:intro}

Mechanistic interpretability assigns causal roles to model components via interchange interventions---swapping internal representations between paired inputs \citep{geiger2021causal, vig2020investigating}. Successful experiments are labeled as ``the mechanism'' for a capability: induction heads for in-context learning \citep{olsson2022context}, the IOI circuit \citep{wang2023interpretability}, the ``lookback mechanism'' for belief tracking \citep{prakash2025lookback}. These labels carry validity assumptions---construct, measurement, internal, external, interpretive---that are rarely tested.

We audit \citet{prakash2025lookback}, who claim that large language models implement belief tracking via three ``lookback'' sub-mechanisms---binding (layers 33--38), answer (layers 52--56), and visibility (layers 10--31)---collectively described as ``pointer dereference via attention.'' We apply 35 validity criteria and identify six untested assumptions.\footnote{We draw on three deposited frameworks: Mechanistic Validity \citep{tower2026mechval}, which provides the 35 criteria; Mechanistic Views \citep{tower2026mechviews}, which asks what \emph{kind} of claim a given method can support; and Mechanistic Reference \citep{tower2026mechref}, which asks whether a mechanism's name refers to anything the method can actually detect. Each is independently deposited; the audit's findings stand without them.} The claim rests on interchange interventions on 80 synthetic stories in Llama-3-70B-Instruct, pre-filtered to cases the model answers correctly, with qualitative replication in two other models. The three sub-mechanisms operate on different token types at different layer ranges, with per-layer subspace ranks from 3 to 19 singular vectors for the pointer concept and up to 320 for visibility (selected from a truncated SVD of the 8192-dimensional residual stream).\footnote{Ranks extracted from released results at \url{https://github.com/Nix07/belief_tracking} (commit \texttt{0579347e}); the SVD uses 500 components. See \texttt{reference/extracted\_results\_llama70b.json} in our repository.}

We produce a verdict on the five-tier scale (Proposed, Causally Suggestive, Mechanistically Supported, Triangulated, Validated) and pre-register six experiments to test the gaps. Section~\ref{sec:discussion} narrows the claim to what the evidence licenses.


\section{Validity Assessment}
\label{sec:validity}

Table~\ref{tab:validity} and Figure~\ref{fig:validity_grid} summarize the assessment. We discuss the most consequential gaps below.

\subsection{Construct validity}

\paragraph{No discriminant validity (C4).}
No task where the lookback pattern should be absent is tested. Candidate negative controls: factual recall without belief tracking, simple copying, single-character stories, non-ToM multi-step retrieval. If the pattern appears on all of these, it is attention-mediated retrieval, not belief tracking.

\paragraph{No convergent evidence (C3).}
All evidence comes from one method family: interchange interventions on residual stream vectors via NNsight. Subspace identification uses the same framework (SVD of intervention-site activations, then learning a binary mask over singular vectors to maximize IIA). No weight analysis, no attention pattern analysis, no independent probing.

The paper's evidence occupies a single level \citep{tower2026mechviews}: the role level, where mechanisms are functional roles validated by intervention success. The subspace identification (DCM) parameterizes the search at a finer grain---selecting directions on the Grassmannian $\text{Gr}(k, d)$---but validates only by IIA, which is a role-level criterion (does the interchange produce the predicted output?). Confirming a genuine subspace claim requires convergent evidence from weight-space analysis: which attention heads implement the retrieval, whether the identified subspace is stable across random seeds, and whether it occupies a consistent location on the Grassmannian across models. The paper provides none of these.

Eleven experiments across three models and dozens of layer-concept combinations look like massive convergent evidence. They are one measurement family applied repeatedly. Quantity of replications within one method does not substitute for an independent method.

\paragraph{Construct import.}
``Belief tracking'' imports cognitive-science vocabulary---Theory of Mind \citep{premack1978chimpanzee}, false-belief attribution \citep{wimmer1983beliefs}---without validating the transfer. The model passes a behavioral ToM test; the internal mechanism gets cognitive-science labels. Behavioral equivalence does not entail mechanistic equivalence.

\subsection{Measurement validity}

\paragraph{No random subspace baseline (M2).}
SVD followed by L1-regularized binary mask selection produces per-layer subspaces of varying rank---typically 3 to 19 singular vectors for the pointer concept, from a truncated SVD of the 8192-dimensional residual stream. No random subspace of matched rank is tested for comparable IIA. In several experiments, subspace IIA approaches full residual stream IIA (e.g., 0.925 vs.\ 1.0 at layer 38 for the pointer concept), suggesting the subspace captures most task-relevant signal. The visibility lookback is particularly striking: the selected ``subspace'' uses 160--320 of 500 SVD components (32--64\% of the truncated basis, 2--4\% of the full residual stream), where a weaker L1 penalty ($\lambda = 0.03$ vs.\ $0.1$) lets the optimizer retain most candidate dimensions. At early layers, the selected subspace achieves higher IIA than the full 500-component basis (e.g., 0.90 vs.\ 0.175 at layer 7), suggesting remaining SVD components destroy counterfactual behavior---consistent with meaningful directional selectivity or optimization artifacts.

\paragraph{Method vacuousness (M2, continued).}
\citet{sutter2025vacuous} demonstrated that unconstrained DAS achieves 100\% IIA on randomly initialized models: a method that succeeds on every system, including random ones, has zero discriminative power. DCM constrains the search to binary masks over a fixed SVD basis rather than arbitrary rotations, so it is more constrained than unconstrained DAS. Whether these constraints are sufficient to avoid vacuousness is an empirical question: if DCM finds low-rank subspaces with high IIA on a randomly initialized model, the method inherits the Sutter problem despite its constraints. We include this as a control in Experiment~2.

\paragraph{No effect size calibration (M1).}
IIA is the sole metric, reported without confidence intervals or statistical tests. For $n = 80$, the standard error on IIA is at most ${\sim}5.6$ percentage points ($\sqrt{0.5 \cdot 0.5 / 80}$, the maximum at $p = 0.5$), placing values below ${\sim}0.60$ within noise of chance for a binary outcome.

\paragraph{Selection on the dependent variable (M3).}
All CausalToM examples were pre-filtered to cases the model answers correctly on both clean and counterfactual prompts. The analysis conditions on the output, then attributes a mechanism to the internal process. Failure cases are excluded. The codebase does generate a train/eval split (80 stories each; \texttt{scripts/patching\_scripts/run\_patching\_exp\_utils.py:465}): the binary mask is optimized on training stories and IIA is reported on held-out validation stories. However, the SVD basis---which defines the candidate directions the mask selects from---is pre-computed offline (\texttt{svd/} directory, not committed to the repository), with no documentation of a held-out procedure. The evaluation is therefore out-of-sample relative to the mask but likely in-sample relative to the feature space.

\subsection{Internal validity}

\paragraph{Full-vector confound (I5).}
Interchange interventions swap the entire residual stream vector (or a large subspace), transferring all information at that position, including the hypothesized ``pointer'' along with everything else. Two interpretations are consistent: (a) the model uses a specific pointer structure and the swap transfers it, or (b) the model uses any information in the residual stream and the swap includes it. Interpretation (b) makes no mechanistic claim beyond ``information flows through the residual stream.''

\paragraph{No epistatic interaction test (I6).}
The three sub-mechanisms are tested independently. Whether disrupting binding changes the answer lookback's behavior beyond the additive prediction is untested. Non-additive interaction would support a single staged mechanism; additive independence would support three separate processes sharing a label.

\paragraph{No rescue reversibility (I7).}
The paper ablates the mechanism (showing necessity) but does not restore a corrupted component to test whether behavior recovers. Rescue experiments distinguish a causal bottleneck from a correlated side-effect.

\subsection{External validity}

\paragraph{Single template (E2).}
All 80 stories follow the CausalToM template with fixed token positions, exactly two characters, and exactly two objects. No structural variation is tested.

\paragraph{No cross-task transfer (E3).}
The paper tests model generality (Gemma-3-27B, Qwen-2.5-14B) without task generality: different story formats, naturalistic text, multi-step reasoning with belief tracking as one component.

\paragraph{No graded response (E5).}
No partial ablation of the subspace is tested for proportional degradation. A dose-response curve would distinguish a sharply localized mechanism from diffuse information storage.

\paragraph{No novel prediction (E6).}
The mechanism is characterized post hoc, with no new behavioral prediction derived from the mechanistic account and tested on held-out data.

\subsection{Interpretive validity}

\paragraph{Referential inflation (V1, V4).}
``Pointer dereference'' imports specificity from programming languages: allocate a reference, store an address, follow the reference, retrieve a value. The evidence shows attention-mediated information flow between token positions. Consider three statements of the same result:

\begin{quote}
\emph{What the experiment shows.} Replacing the residual stream vector at the final token position with the vector from a counterfactual story changes the model's answer to match the counterfactual, for a 19-dimensional subspace of a 500-component SVD basis, at layers 52--56, on 80 template-matched stories.

\emph{What this licenses.} Information at this position and layer range is causally necessary for producing the belief-consistent answer.

\emph{What the paper claims.} The model implements pointer dereference: it stores an address, follows the reference, and retrieves the payload.
\end{quote}

The gap between the second and third statement is not supplied by any experiment---it is supplied by the vocabulary. Substituting ``pointer'' $\rightarrow$ query vector, ``address'' $\rightarrow$ key vector, ``payload'' $\rightarrow$ value content, ``dereference'' $\rightarrow$ attention, the finding reduces to: the model uses attention to retrieve information from earlier tokens. This is what attention does by construction. \citet{tower2026mechref} calls this pattern \emph{referential inflation}: the labels carry referential content---computational specificity---that the method cannot detect. Strip the metaphor and the level jump becomes visible.

\paragraph{No falsification target.}
No result would have falsified the lookback hypothesis as stated. High IIA supports it, low IIA means ``the mechanism does not operate here,'' and matched full-vector and subspace IIA means ``the subspace captures the information.'' Every outcome confirms.

\paragraph{The third-answer result.}
The paper's strongest evidence is the answer pointer experiment (Figure 4b): swapping the ``pointer'' while preserving the ``payload'' produces a third answer~$C$ distinct from both the original answer~$A$ and the counterfactual answer~$B$. This demonstrates decomposable internal structure---a random-noise perturbation would degrade $A$, not produce a coherent alternative.

The structure, however, is predicted by the attention mechanism itself. QK circuits determine \emph{where} to attend (the pointer); OV circuits determine \emph{what} to retrieve (the payload). Swapping directions that alter QK routing makes the model attend to a different earlier position and retrieve whatever content is stored there via OV---producing a coherent third answer corresponding to the content at the redirected position. Every transformer attention head exhibits this pointer-payload decomposition by construction; the question is whether the specific subspace identified by DCM captures belief-specific routing or generic QK structure.

We confirm the decomposable structure as genuine and test the belief-specificity question in Experiments~2 and~3 (Section~\ref{sec:experiments}). Experiment~2 includes a specific test: if random subspaces of matched rank also produce coherent third answers at above-chance rates, the decomposition is a property of the SVD basis geometry, not of the identified directions. We define ``coherent'' quantitatively before seeing the data: a third answer counts as coherent if it is a token from the story's entity vocabulary (character names, objects, locations) distinct from both $A$ and $B$. The coherent-third-answer rate on the identified subspace must exceed the rate on random subspaces ($p < 0.01$, rank test) to support directional specificity.

\subsection{Summary}

\begin{table}[t]
\centering
\caption{Validity assessment against 35 criteria \citep{tower2026mechval}. Two criteria fail outright (V1: level declaration, V4: anthropomorphism check). Criteria not discussed individually are rated in the supplementary audit document. The four additional criteria beyond the original 31 (I3, I5, I8, I11) were added in framework version 11; see Amendment~10a.}
\label{tab:validity}
\begin{tabular}{lcccc}
\toprule
Category & Confirmed & Partial & Not met & Fails \\
\midrule
Construct (C1--C5) & 1 & 2 & 2 & 0 \\
Measurement (M1--M7) & 0 & 2 & 5 & 0 \\
Internal (I1--I11) & 1 & 2 & 8 & 0 \\
External (E1--E6) & 0 & 2 & 4 & 0 \\
Interpretive (V1--V5) & 0 & 2 & 1 & 2 \\
\midrule
\textbf{Total (35)} & \textbf{2} & \textbf{10} & \textbf{21} & \textbf{2} \\
\bottomrule
\end{tabular}
\end{table}

\begin{figure}[t]
\centering
\includegraphics[width=0.85\textwidth]{figures/fig6_validity_grid.pdf}
\caption{Validity assessment across 35 criteria grouped into five categories. Most criteria are ``not met'' (21/35), reflecting the paper's minimal reporting of measurement and internal validity evidence. Two criteria fail outright (V1: level declaration, V4: anthropomorphism check). The ten partial ratings correspond to criteria where the paper provides some relevant evidence without meeting the full standard.}
\label{fig:validity_grid}
\end{figure}

This is Causally Suggestive on the five-tier scale. The paper establishes necessity (I1: ablation degrades performance); the bottleneck to Mechanistically Supported is independent-method evidence (C3), negative controls (C4), and a random subspace baseline (M2).


\section{The Unity Problem}
\label{sec:unity}

The paper treats three ``lookback mechanisms'' as instances of a single computation. They operate at non-overlapping layer ranges, on different token types, with per-layer subspace ranks that differ by an order of magnitude (3--10 for binding vs.\ 160--320 for visibility; Table~\ref{tab:unity}).

\begin{table}[t]
\centering
\caption{The three ``lookback mechanisms'' differ in every measurable property. Subspace ranks are per-layer (selected via L1-regularized binary mask over SVD components) and vary across layers within each mechanism. Ranks extracted from released results at commit \texttt{0579347e}; exact values depend on layer and L1 weight ($\lambda = 0.1$ for binding/answer, $\lambda = 0.03$ for visibility).}
\label{tab:unity}
\begin{tabular}{lccc}
\toprule
Lookback & Layers & Tokens & Per-layer rank \\
\midrule
Binding & 33--38 & State tokens & 3--10 \\
Answer & 52--56 & Final token (``:'') & 18--19 \\
Visibility & 10--31 & Visibility sentence & 160--320 \\
\bottomrule
\end{tabular}
\end{table}

The unity claim rests on shared vocabulary: all three involve ``lookback'' (attention to an earlier token) and ``dereference'' (retrieval at the attended position). This describes the attention mechanism itself---every QK circuit attending to a previous position and retrieving content through OV fits the definition.

The critical test is coupling: if binding (layers 33--38) is disrupted, does the answer lookback (layers 52--56) degrade? Sequential dependence would indicate one staged mechanism. Independent disruptability would indicate separate processes sharing a naming convention. The paper does not perform this test.


\section{Cognitive Science Context}
\label{sec:cogsci}

The paper frames its contribution as understanding ``how LMs track beliefs,'' importing Theory of Mind from cognitive science. Three lines of evidence bear on this.

Vegner et al.\ (\citeyear{vegner2025systematicity}) argue that the field conflates behavioral systematicity (correct outputs across variations) with representational systematicity (compositionally structured internal representations). The Lookback paper shows the former and infers the latter. CausalToM uses exactly two characters, first-order beliefs only, and a single story template. Representational systematicity predicts clean extension to three or more characters, nested beliefs, and structurally novel stories. Behavioral matching does not require any of these.

The single-mechanism assumption is itself risky. Apperly and Butterfill (\citeyear{apperly2009humans}) proposed that human belief tracking involves two distinct systems. Phillips et al.\ (\citeyear{phillips2015second}) failed to replicate the strongest evidence for a single automatic system \citep{kovacs2010social}, finding the original effects better explained by timing artifacts. If cognitive scientists cannot establish that humans have a single unified belief-tracking mechanism, claiming language models have one is premature. The paper's own data reinforce this concern: three dissociable sub-effects at different layers, tokens, and dimensions.

Human belief attribution is content-selective: people track false beliefs about presence differently than about absence \citep{schuwerk2015implicit}. CausalToM tests one content type. A ``pervasive'' mechanism should handle beliefs about location, existence, intentions, and knowledge states. Ruis et al.\ (\citeyear{ruis2025reevaluating}) make the general point: passing a simplified ToM behavioral test and finding a correlated internal pattern does not establish that the pattern implements ToM in any sense that transfers from cognitive science.


\section{Pre-Registered Experiments}
\label{sec:experiments}

An audit of another group's work can easily become a post-hoc narrative: run many experiments, report the ones that tell a story. We pre-register six experiments and commit predictions before seeing data so the reader can verify which outcomes we expected and which surprised us. Three experiments are primary (run first), three secondary (run if primary experiments succeed). All use Llama-3-70B-Instruct via NNsight, reproducing the original setup. All $p$-values are Benjamini-Hochberg corrected across the six experiments. The pre-registration was committed at tag \texttt{prereg-amendment-5} before results were computed; subsequent amendments (7--9) are SHA-tagged in \texttt{AMENDMENTS.md}.

\subsection{Held-out basis validation (M2, M3)}

The original codebase validates the binary mask out-of-sample (mask trained on 80 stories, IIA evaluated on a separate 80). The SVD basis itself, however, is pre-computed offline with no documented held-out procedure. This experiment tests whether the reported subspaces generalize when the SVD basis is also computed on disjoint data. Protocol: partition a larger pool of CausalToM stories (320+, generated from the same templates) into three non-overlapping sets---SVD estimation, mask training, and IIA evaluation. Recompute the SVD on the estimation set, fit the binary mask on the training set using the original hyperparameters, and evaluate IIA on the held-out evaluation set. Repeat for each layer $\times$ concept combination.

\paragraph{Pre-committed interpretations.} If IIA on held-out data matches the released values within $\pm 0.05$ (4 examples on an 80-story evaluation set), the SVD-basis-leakage concern is closed and we report M3 as satisfied for the feature space. If IIA drops substantially (below 0.5 or below the 95th percentile of the random subspace distribution from Experiment~2), the reported subspaces are artifacts of basis overfitting and the claim requires recomputed SVD with proper held-out evaluation.

\subsection{Matched-capacity controls (M2)}

Two controls with matched optimization budget, evaluated on the same held-out set from Experiment~1:

\paragraph{Random subspace baseline.} For each layer's identified subspace of rank $r$, sample 200 random $r$-of-500 index sets from the same SVD basis (uniformly without replacement, matching the paper's binary-mask method). Compute IIA for each on the held-out evaluation set. The identified subspace should rank in the top 1\% ($p < 0.01$, one-sided rank test with corrected $p = (k + 1)/(n + 1)$). This is a floor control: the identified subspaces are optimized to maximize IIA while random selections are not, so the identified subspace beating random draws is expected. Additionally, we record the coherent-third-answer rate (see Section~\ref{sec:validity}) for each random subspace to test whether the decomposable structure is direction-specific.

\paragraph{Shuffled-label control.} For two primary layer-concept combinations (one binding: \texttt{binding\_lookback/address\_and\_payload} at layer 34; one answer: \texttt{answer\_lookback/pointer} at layer 52), re-run the SVD + mask selection procedure with identical hyperparameters (Adam, $\text{lr} = 0.1$, L1 weight matching the original) on 100 permutations of the belief labels---shuffling which story's belief state serves as source for interchange intervention. Crucially, shuffled masks are \emph{trained} on shuffled assignments but \emph{evaluated} on real counterfactual targets (held-out set). The real-label subspace should beat the shuffled-label distribution ($p < 0.01$, rank test). The remaining four layer-concept combinations are tested as exploratory with 20 permutations each ($p$-floor $= 0.048$).

\paragraph{Vacuousness control.} Apply the DCM procedure (SVD of activations, binary mask optimization with identical hyperparameters) to a randomly initialized Llama-3-70B. If DCM finds low-rank subspaces with IIA exceeding chance, the method inherits the \citet{sutter2025vacuous} vacuousness problem despite its constraints, and the subspace identification on the trained model is uninformative.

\paragraph{Pre-committed interpretations.} If real-label subspaces beat both random and shuffled-label baselines on held-out data, M2 is satisfied and the belief-specificity concern narrows to the task-generality question (C4, Experiment~3). If the vacuousness control produces high-IIA subspaces on the random model, DCM lacks referential content and the subspace results are methodologically void regardless of other outcomes. If shuffled-label subspaces achieve comparable IIA, the mask selection captures generic high-variance directions rather than belief structure, and the claim requires revision. If both shuffled and random subspaces approach the identified subspace's IIA, the subspace identification adds no information beyond dimensionality matching. Additionally, for the third-answer test: if random subspaces of matched rank produce coherent third answers (story-vocabulary tokens distinct from both original and counterfactual answers) at rates comparable to the identified subspace, the decomposable structure is a property of the basis geometry rather than of the identified directions.

\subsection{Task-specificity controls (C4)}

Run the interchange intervention protocol on programmatically generated stories (120 per condition) across four conditions that isolate binding and belief attribution:

\begin{enumerate}
\item Echo/copy (no binding, no belief): ``Alice said: `the X is in loc1.' What did Alice say?''---pure retrieval, the strong control.
\item Entity tracking (binding, no belief): ``Alice put X in loc1. Bob moved X to loc2. Where is X?''---a weak control, as the authors' own work predicts the pattern may appear here.
\item True-belief matched stories (binding, minimal belief): same CausalToM structure but the character's belief matches reality.
\item First-order false belief: standard CausalToM (the original task).
\end{enumerate}

Use the same layer ranges, subspace dimensions, and IIA metric across all conditions.

\paragraph{Pre-committed interpretations.} If the lookback pattern is absent on echo/copy and entity tracking but present on false belief, C4 is satisfied: the subspace captures belief-specific structure. If the pattern appears on entity tracking but not echo/copy, the mechanism is entity-state retrieval rather than belief tracking---a substantive correction of the interpretive claim, consistent with the authors' own entity-tracking prior work. If the pattern appears on all conditions including echo/copy, it is generic attention-mediated retrieval, not a belief-tracking mechanism.

\subsection{Compositional generalization (C3, C5)}
\label{sec:composition}

\citet{vegner2025systematicity} distinguish behavioral systematicity (correct outputs across variations) from representational systematicity (compositionally structured internal representations). The Lookback paper shows the former and infers the latter. CausalToM uses exactly two characters, first-order beliefs, and a single story template. If the lookback mechanism is a genuine compositional representation of belief, it should extend to structurally novel cases. If it is a template-matched retrieval circuit, it will not.

We test three axes of generalization (120 stories per condition):

\begin{enumerate}
\item \textbf{Character scaling.} CausalToM variants with 3, 4, and 5 characters. The binding lookback assigns ordinal indices (OI-1, OI-2) to two characters; a compositional scheme should extend to OI-3, OI-4, OI-5 with additional binding slots. A template-matched circuit will fail or degrade.
\item \textbf{Nested beliefs.} Second-order false belief: ``What does Alice think Bob believes Object2 contains?'' The binding lookback should extend to represent the nesting structure. If the same subspace handles first- and second-order beliefs at the same layers and ranks, the mechanism generalizes. If second-order requires different layers or substantially higher rank, the mechanism is first-order-specific.
\item \textbf{Content selectivity.} Human belief tracking is content-selective: people track false beliefs about presence differently than about absence \citep{schuwerk2015implicit}. We test beliefs about object location, object existence, character intentions, and character knowledge states. A unified mechanism should handle all content types with comparable IIA at comparable layers. Selective failure indicates a family of content-specific circuits rather than one mechanism.
\end{enumerate}

For each condition, we apply the full DCM procedure (SVD + binary mask optimization with the original hyperparameters) and compare: (a) whether the identified subspace overlaps with the original CausalToM subspace (cosine similarity between mask vectors), (b) whether IIA is comparable, and (c) whether the active layer ranges shift.

\paragraph{Pre-committed interpretations.} If the lookback pattern generalizes across all three axes with comparable IIA, layer ranges, and subspace dimensions, the representational systematicity claim is supported and ``belief tracking'' is a defensible label. If the pattern degrades on character scaling ($\text{IIA} < 0.5$ at 3+ characters), it is a two-character retrieval circuit. If it fails on nested beliefs, it is first-order-specific. If it shows content selectivity (different layers or subspaces for different belief types), ``lookback'' is a family name for multiple circuits rather than one mechanism---paralleling the dissociation observed in human belief attribution \citep{apperly2009humans}. Any of these outcomes is informative: degradation patterns reveal the mechanism's actual scope.

\subsection{Failure case analysis (I1)}
\label{sec:failure}

Run identical interventions on CausalToM stories where Llama-3-70B-Instruct answers incorrectly. The original paper excludes these by design. To ensure sufficient failures for a powered test ($n_{\text{fail}} \geq 20$), we generate a pool of 500+ CausalToM stories rather than using the original 80.

\paragraph{Pre-committed interpretations.} If lookback IIA is significantly lower on incorrect-answer cases ($p < 0.01$, permutation test), I1 is satisfied: the mechanism is absent when the model fails, supporting a causal role. If the lookback is present and intact on failure cases, it is a correlate of processing rather than a cause of correct answers. If the model's accuracy yields fewer than 20 failures even on 500 stories, we report the experiment as underpowered rather than drawing a conclusion.

\subsection{Cross-stage mediation (I6)}
\label{sec:mediation}

Test whether the binding and answer sub-mechanisms form a causal chain or operate independently, replacing the super-additivity test originally proposed.\footnote{Super-additive behavioral loss is ambiguous: independent stages in a serial pipeline can produce super-additive loss through off-distribution compounding, and components of one mechanism can appear additive.} Protocol: (a) intervene on the binding representation at layers 33--38; (b) measure the resulting change in the answer-pointer representation at layers 52--56; (c) intervene on the answer-pointer representation; (d) test whether the answer-pointer intervention mediates the binding intervention's effect on output. All interventions use mean-ablation (replacing the subspace component with its dataset mean) to keep activations on the data manifold.

\paragraph{Pre-committed interpretations.} If binding interventions propagate through the answer-pointer representation and the answer-pointer mediates binding's effect on output, the unity claim is supported: binding $\to$ pointer $\to$ answer forms a causal chain. If the stages independently affect output without mediating each other, they are separate processes sharing a label, and ``lookback mechanism'' as a unified construct is not warranted. If the mediation is partial, the truth is intermediate and we report the mediation proportion.


\section{Results}
\label{sec:results}

All experiments use Llama-3.1-70B-Instruct via NNsight/NDIF remote execution, applying the paper's released SVD bases and singular-vector masks to compute interchange intervention accuracy (IIA).

\subsection{Measurement validity: replication and baselines (E1--E2)}

\paragraph{E1: Positive control replication.}
We replicate the paper's interchange interventions on Llama-3.1-70B-Instruct (paper uses 3.0). L34 binding replicates within $\pm 0.05$ (IIA $= 0.988$ vs.\ paper's $0.963$). The answer subspaces show \emph{higher} IIA than reported: L38 $= 1.0$ (paper: $0.925$), L52 $= 1.0$ (paper: $0.775$), L53 $= 1.0$ (paper: $0.550$). L36 binding drops ($0.225$ vs.\ paper's $0.750$). L52 and L53's upward shifts suggest seed 10 underestimates typical IIA; multi-seed replication confirmed L53 variance spanning $0.55$--$0.99$ across seeds (Figure~\ref{fig:replication}).

\begin{figure}[t]
\centering
\includegraphics[width=0.85\textwidth]{figures/fig4_replication.pdf}
\caption{Positive control replication (E1). IIA for all six subspaces: Prakash et al.\ on Llama-3.0-70B-Instruct (grey) vs.\ this work on Llama-3.1-70B-Instruct (blue). Four subspaces replicate within $\pm 0.13$; L36 fails to replicate (paper $0.75$, ours $0.23$); L52 and L53 increase to ceiling. Numbers above bars show the delta.}
\label{fig:replication}
\end{figure}

\paragraph{E2: Shuffled-label and random-subspace controls.}
All six subspaces beat random subspaces of matched rank ($p = 0.005$, rank test, $n = 200$ random draws). Against shuffled-label controls ($n = 20$ permutations), four pass at $p < 0.05$: L34 binding ($0.010$), L52 answer ($0.010$), L38 answer ($0.048$), L35/L36 binding ($0.048$ each). L53 does not reach significance ($p = 0.286$), consistent with high seed-to-seed variance. Real-label IIA ($0.55$--$0.96$) is well separated from shuffled-label means ($0.47$--$0.50$).

\subsection{Construct validity: specificity and generalization (E10--E11)}

\paragraph{E10: Cross-template generalization.}
Answer subspaces trained on template 2 (the paper's default) transfer to template 0 ($n = 23$): L38 IIA $= 0.957$, L52/L53 $= 0.783$. Not template-specific.

\paragraph{E11: Cross-task transfer.}
On factual recall ($n = 34$), L38 achieves IIA $= 0.529$ while L52 and L53 both achieve $0.0$. On color-property retrieval ($n = 6$), L38 achieves $0.333$; L52/L53 again achieve $0.0$. Arithmetic ($n = 0$) yielded no viable pairs. L38 partially transfers to non-belief retrieval tasks; L52 and L53 do not. This dissociation suggests L38 encodes a more generic entity-retrieval mechanism that the belief circuit routes through, while L52/L53 are task-specific (Figure~\ref{fig:crosstask}).

\begin{figure}[t]
\centering
\includegraphics[width=\textwidth]{figures/fig5_crosstask.pdf}
\caption{Cross-task dissociation between L38 and L52. L38 partially transfers to factual recall (IIA $= 0.53$) and action recall (IIA $= 0.89$), while L52 drops to zero on non-belief tasks. Both subspaces achieve near-ceiling IIA on belief-relevant tasks (belief framing, false belief, observability). The gap between the two curves quantifies how much L38 acts as a generic retrieval mechanism compared to L52's task-specific function.}
\label{fig:crosstask}
\end{figure}

\subsection{Belief-vs-binding dissociation (E12--E15)}
\label{sec:dissociation_results}

Pre-registered in Amendment~9, these experiments test whether the subspaces encode generic entity-state bindings or character-perspective information.

\paragraph{E12: Question-framing control ($n = 35$ shared pairs).}
We generate clean/counterfactual pairs from the CausalToM template and reframe the question four ways, preserving the story and answer: \emph{belief} (``What does X believe the container contains?''), \emph{action recall} (``What did X put in the container?''), \emph{reality state} (``What is inside the container?''), \emph{fill completion} (``X filled the container with'').

Results in Table~\ref{tab:framing}. The subspaces are framing-\emph{sensitive}: belief framing produces IIA $= 1.0$ across all subspaces, action recall yields partial transfer (L38 $= 0.886$, L52 $= 0.400$, L53 $= 0.057$), reality-state and fill-completion framings produce IIA $= 0.0$. This \emph{overturns} our pre-registered prediction that all framings would be within $0.10$ of the belief baseline. Pre-committed interpretation: ``the subspace may be sensitive to epistemic question framing.''

If the subspaces encoded generic entity-state bindings (``what drink is in what container''), reality-state and fill-completion questions---which request the same underlying fact---should activate the same subspace. They do not. The subspaces encode information relevant specifically to character-perspective queries.

\begin{table}[t]
\centering
\caption{E12: IIA by question framing ($n = 35$ shared pairs) with 95\% Wilson confidence intervals. The subspaces specifically respond to belief framing, partially to action recall, and not at all to perspective-free queries about reality or fill completion. See Figure~\ref{fig:framing}.}
\label{tab:framing}
\begin{tabular}{lccc}
\toprule
\textbf{Framing} & \textbf{L38} & \textbf{L52} & \textbf{L53} \\
\midrule
Belief (``What does X believe\ldots'') & $1.00\ [.90, 1.0]$ & $1.00\ [.90, 1.0]$ & $1.00\ [.90, 1.0]$ \\
Action recall (``What did X put\ldots'') & $.89\ [.74, .96]$ & $.40\ [.25, .57]$ & $.06\ [.01, .19]$ \\
Reality state (``What is inside\ldots'') & $.00\ [.00, .10]$ & $.00\ [.00, .10]$ & $.00\ [.00, .10]$ \\
Fill completion (``X filled\ldots with'') & $.00\ [.00, .10]$ & $.00\ [.00, .10]$ & $.00\ [.00, .10]$ \\
\bottomrule
\end{tabular}
\end{table}

\begin{figure}[t]
\centering
\includegraphics[width=\textwidth]{figures/fig2_framing_dissociation.pdf}
\caption{Question framing dissociation (E12, $n = 35$). The same clean/counterfactual story pairs produce IIA $= 1.0$ under belief framing, graded transfer under action recall (L38 $= 0.89$, L52 $= 0.40$), and IIA $= 0.0$ under reality-state or fill-completion framing. Error bars show 95\% Wilson confidence intervals. The underlying entity-state fact is identical across all four framings; the subspace responds to the epistemic framing of the question. This dissociation is the strongest evidence that the subspaces encode character-perspective information.}
\label{fig:framing}
\end{figure}

\paragraph{E13: Distractor insertion ($n = 80$ baseline, $n = 21$ early, $n = 0$ mid/late).}
We insert irrelevant drink mentions at three positions: early (after scene-setting), mid (between fill actions), late (before the question). Baseline: L38/L52 IIA $= 1.0$, L53 $= 0.7$. Mid- and late-position distractors destroy behavioral accuracy entirely (zero surviving pairs). Early distractors retain 21 pairs, but IIA drops to $0.0$ across all subspaces. The subspace is positionally fragile: even minor template perturbations disrupt the mechanism. This limits the belief-tracking claim to the template structure used in training, though it does not bear on belief vs.\ generic retrieval.

\paragraph{E14: True false-belief test ($n = 80$, 100\% behavioral accuracy).}
A character places a drink, a second character swaps it unobserved, and the question asks what the first character \emph{believes} the container holds. The model answers all 80 correctly (reporting the original drink, not the swapped one).

L38 and L52 both achieve IIA $= 0.975$ (95\% Wilson CI $[0.913, 0.993]$; L53 $= 0.388$). This \emph{overturns} our prediction of IIA $< 0.30$. Pre-committed interpretation: ``the subspace genuinely tracks the character's epistemic state even when it diverges from reality.'' Subspaces trained where belief equals reality generalize to stories where they diverge.

The reality-question control (asking what is \emph{actually} in the container) yielded zero viable pairs---the model could not reliably answer these. Uninformative.

\paragraph{E15: Observability dissociation ($n = 80$, 88\% behavioral accuracy).}
Template 1 creates an observability asymmetry: character 1 can observe character 2's actions, not vice versa. We ask character 1 about character 2's container (observed-other condition). The model answers 212 of 240 generated pairs correctly (88\%); we evaluate the first 80 passing pairs.

L38 IIA $= 0.988$ (95\% Wilson CI $[0.932, 0.998]$), L52 $= 0.975$ $[0.913, 0.993]$ (L53 $= 0.213$). Again overturns our prediction of IIA $< 0.30$. The subspaces generalize to knowledge acquired through observation, not just through the character's own actions.

The self-container control yielded zero viable pairs. Uninformative.

\subsection{Cross-model dependency structure (X2)}
\label{sec:crossmodel}

Prakash et al.'s cross-model appendix tests Qwen2.5-14B-Instruct and Llama-3.1-405B-Instruct with full-residual interchange interventions only, acknowledging that ``subspace interchange intervention experiments were not conducted'' yet concluding that these models ``employ the same underlying mechanism.'' Full-residual patching replaces all information at a position simultaneously and cannot distinguish address from payload.

We apply the paper's own three-condition mismatch test---their strongest causal design---to Qwen2.5-14B-Instruct ($n = 200$ pairs, 12 layers per lookback type). This test patches recalled tokens only (address updated, pointer stale), lookback tokens only (pointer updated, address stale), or both (coherent). The mismatch signature---singles fail, pair succeeds---demonstrates dependency structure without requiring subspace fitting.

\paragraph{Answer lookback: mismatch signature at layers 32--38.}
Recalled-only IIA $= 0.005$ $[0.001, 0.028]$, lookback-only IIA $= 0.000$ $[0.000, 0.019]$, both IIA $= 1.000$ $[0.981, 1.0]$ (95\% Wilson CIs, L32). The signature holds from L32 through L38 (67--79\% of Qwen's 48 layers). Layer 30 shows a transition (both $= 0.440$, singles near zero). The proportional depth is comparable to Llama-3-70B's answer range (layers 52--56, 65--70\% of 80 layers). Patching either component alone has no effect; patching both restores full counterfactual transfer.

\paragraph{Binding lookback: anti-mismatch.}
Recalled-only IIA $= 1.0$ across all layers (patching state tokens alone suffices). Lookback-only IIA rises from $0.0$ at L12 to $1.0$ by L24. From L24 onward, each component alone succeeds but the combined patch yields IIA $= 0.0$. The binding mechanism on Qwen does not require coordination between recalled and lookback positions---the combined patch creates an internally incoherent state. This is a genuine structural difference that full-residual patching cannot detect.

The answer lookback dependency structure transfers to Qwen; the binding structure does not. ``Same underlying mechanism'' is accurate for answer lookback and false for binding---partial transport that the paper's binary claim cannot express.

\subsection{Rescue reversibility (I7, L38, $n = 200$)}

We zero-ablate the L38 answer subspace (rank 3) and measure five conditions: (A) clean baseline, (B) zero-ablated, (C) ablated then restored from original activation, (D) ablated then restored from a different sample, (E) ablated then restored with a random subspace of matched rank. Clean accuracy $= 1.0$; ablated $= 0.92$; rescued $= 1.0$; wrong-rescue $= 0.545$; random-rescue $= 0.93$. The ablation drop (8\%) is small---the model is robust to subspace removal, consistent with redundant pathways in the 8192-dimensional residual stream. The wrong-rescue drop to $0.545$ confirms the subspace carries sample-specific information: restoring the wrong sample's projection actively degrades performance below the ablated baseline. Random-subspace restoration ($0.93$) does not help, confirming directional specificity. The remaining answer subspaces (L52, L53) were not tested due to NDIF instability.

\subsection{Experiments with partial or incomplete results}

\paragraph{E12v2: Question-framing replication ($n = 250$ shared pairs, partial).}
We rerun the question-framing control at $n = 250$ with per-framing accuracy filtering. Belief framing confirms the $n = 35$ result: IIA $= 1.0$ for all three subspaces (L38 $[0.985, 1.0]$, L52 $[0.985, 1.0]$, L53 $[0.985, 1.0]$). Action recall shows IIA near ceiling (L38 $= 0.98$, L52 $= 1.0$, L53 $= 1.0$), higher than the $n = 35$ result (L38 $= 0.886$, L52 $= 0.400$, L53 $= 0.057$). This shift is attributable to a methodological difference: the $n = 250$ design filters pairs per-framing (requiring model accuracy under each framing before computing IIA), then intersects, whereas the $n = 35$ design filtered on belief framing only. Per-framing filtering selects stories where the model handles action recall well, biasing the action-recall IIA upward. Reality-state and fill-completion framings were not computed due to NDIF session termination. The primary dissociation---belief versus reality/fill---cannot be confirmed or disconfirmed at higher power until those framings complete.

\paragraph{Remaining incomplete experiments.}
E7 (necessity) and E8 (layer spread): filtered pairs cached but no IIA computed. M1 (dose-response graded ablation): filtered pairs cached, no graded results. E9 (rank decomposition): the per-direction intervention produced zero IIA at all ranks including full-rank reconstruction, despite the original subspace achieving IIA $> 0.9$---a methodology failure in the per-direction intervention, not a scientific finding. E4 (surface heuristic) and E5 (compositional generalization): only dry-run or zero-pair results. These experiments do not affect the central conclusions, which rest on E1--E2 (measurement validity), E10--E11 (specificity), E12--E15 (the belief-vs-binding dissociation), and X2 (cross-model dependency structure).

\subsection{Confounding sensitivity (I8, exploratory)}

We compute E-values \citep{vanderweele2017sensitivity} for all 20 positive effects (IIA $>$ baseline) to assess robustness to unmeasured confounders. $E = \text{RR} + \sqrt{\text{RR} \cdot (\text{RR} - 1)}$, where $\text{RR} = \text{IIA} / \text{baseline}$ and baseline $= 0.5$ (chance under random subspace swap).

19 of 20 achieve $E \geq 2.0$ (Figure~\ref{fig:evalues}). The single fragile effect is L38 on factual recall ($E = 1.31$, IIA $= 0.529$), consistent with L38's partial cross-task transfer (E11). Replication, false belief, observability, and belief framing achieve $E = 3.31$--$3.41$. This analysis was added after initial pre-registration (Amendment~10c) and is exploratory.

\begin{figure}[t]
\centering
\includegraphics[width=\textwidth]{figures/fig3_evalues.pdf}
\caption{Confounding sensitivity analysis (I8, exploratory). E-values for all 20 positive effects (IIA $>$ chance), sorted by magnitude. An E-value quantifies how strong an unmeasured confounder would need to be to explain away the observed effect. The dashed line marks $E = 2.0$; effects above this threshold are considered robust. Only one effect falls below: L38 on factual recall ($E = 1.31$), consistent with L38's partial cross-task transfer.}
\label{fig:evalues}
\end{figure}

\subsection{Summary}

Figure~\ref{fig:heatmap} summarizes IIA across all completed experiments.

\begin{figure}[t]
\centering
\includegraphics[width=0.75\textwidth]{figures/fig1_iia_heatmap.pdf}
\caption{IIA across all completed experiments for the three answer-pointer subspaces. L38 and L52 achieve near-ceiling IIA on belief-relevant tasks (replication, false belief, observability, belief framing) and zero on perspective-free queries (reality state, fill completion). L53 is unreliable across conditions, consistent with its high seed-to-seed variance (E1) and failure to reach significance against shuffled labels (E2). Factual recall transfers partially to L38 (0.53) but not L52/L53, supporting the L38-as-generic-retrieval interpretation. All sample sizes are $n = 80$ unless noted.}
\label{fig:heatmap}
\end{figure}

\begin{table}[t]
\centering
\caption{Summary of pre-registered experiments with completed results. Predictions marked with $\dagger$ were overturned by the data; all overturned predictions strengthened rather than weakened the paper's claim. L53 is excluded from the verdict column due to high seed-to-seed variance (E1) and failure to reach significance against shuffled labels (E2).}
\label{tab:summary}
\begin{tabular}{lp{3.2cm}p{3.2cm}l}
\toprule
\textbf{Exp.} & \textbf{Prediction} & \textbf{Result (L38/L52)} & \textbf{Verdict} \\
\midrule
E1 & Replicates within $\pm 0.05$ & L34: $+0.025$; L38/L52: higher & Passes \\
E2 & Beats shuffled/random & $p < 0.05$ (5 of 6) & Passes \\
E10 & Cross-template transfer & L38: $0.957$ & Passes \\
E11 & No cross-task transfer & L38: $0.529$; L52: $0.0$ & Mixed \\
E12$^\dagger$ & Framing-invariant & Framing-sensitive & Overturned \\
X2 & Same mechanism cross-model & Answer transfers, binding differs & Partial \\
I7 & Rescue recovers accuracy & 8\% drop, rescue to 1.0 & Partial \\
E13 & Robust to distractors & Fragile ($n=0$ mid/late) & Fails \\
E14$^\dagger$ & False-belief IIA $< 0.30$ & IIA $= 0.975$ & Overturned \\
E15$^\dagger$ & Observed-other IIA $< 0.30$ & IIA $= 0.988$ & Overturned \\
\bottomrule
\end{tabular}
\end{table}


\section{Discussion}
\label{sec:discussion}

We did not predict this verdict. We expected the subspaces to encode generic entity-state bindings that happen to align with beliefs on the paper's task. Three of four dissociation experiments overturned their pre-registered predictions, each time \emph{strengthening} the paper's claim.

E14 is the strongest result. CausalToM never tests false beliefs: \texttt{Dataset.set\_story()} uses only \texttt{template["context"]}, excluding the \texttt{causal\_event} field, and always queries characters about their own containers (\texttt{set\_character=rc, set\_container=rc}), so belief equals reality for every pair. We constructed 80 stories with genuine drink swaps that the queried character does not observe. The model answers all 80 correctly (reporting the believed drink, not the swapped one), and the subspaces achieve IIA $= 0.975$ at L38 and L52. Subspaces trained where belief equals reality generalize to stories where they diverge---evidence the paper never collected.

E15 extends this. We ask a character about a container filled by someone else, where the queried character observed the filling. L38 IIA $= 0.988$, L52 $= 0.975$. The mechanism generalizes to knowledge acquired through observation, not just through the character's own actions.

E12 is most informative. We predicted framing-invariance: belief, action-recall, reality-state, and fill-completion questions would produce matched IIA, indicating generic entity-state lookup. The opposite occurred. The subspaces activate for character-perspective queries (belief IIA $= 1.0$), show graded transfer to action recall (L38 $= 0.886$, L52 $= 0.400$), and are silent for perspective-free queries (reality-state and fill-completion IIA $= 0.0$). The underlying entity-state fact---which drink is in which container---is identical across all four framings. The subspace responds to the framing, not the fact.

This resolves the central question. The subspaces do not encode ``what entity is associated with what state'' (which would predict framing-invariance). They encode information accessible through character-perspective queries---information that updates when characters observe events (E15) and tracks beliefs diverging from reality (E14). Referential inflation (Section~\ref{sec:validity}) was warranted as an a priori worry, but the mechanism is more specific than generic retrieval.

\paragraph{The L38 vs.\ L52/L53 dissociation.}
L38 partially transfers to factual recall (E11: IIA $= 0.529$) and generalizes across templates (E10: IIA $= 0.957$; Figure~\ref{fig:crosstask}). L52 and L53 show zero transfer to non-belief tasks. L38 also shows the strongest transfer to action-recall framing (E12: $0.886$ vs.\ L52's $0.400$). This suggests L38 encodes a more generic entity-retrieval mechanism that the belief circuit routes through, while L52 is task-specific. The paper's three-stage architecture (binding $\to$ answer $\to$ visibility) may be better characterized as: generic retrieval (L38) feeding into task-specific formatting (L52/L53) feeding into input gating (visibility layers). L53 is unreliable: it fails to reach significance against shuffled labels (E2: $p = 0.286$) and shows high seed-to-seed variance (E1: paper reports $0.55$ at seed 10; we observe $1.0$ at seed 42).

\paragraph{Limitations.}
E13 reveals positional fragility: minor template perturbations destroy behavioral accuracy (mid/late: zero surviving pairs) or subspace function (early: IIA $= 0.0$). The belief-tracking claim holds within the training template structure; the mechanism may not survive naturalistic text. E4--E9 could not be completed due to NDIF server instability and model-accuracy filtering; compositional generalization (E5) remains open. All Llama experiments use one model (Llama-3.1-70B-Instruct, a minor version update from the paper's 3.0). X2 extends to Qwen2.5-14B-Instruct, where the answer lookback dependency structure transfers but the binding mechanism operates differently. We use the paper's released SVD bases and masks rather than recomputing them, so our results test \emph{generalization} to novel stimuli rather than \emph{replicability} from scratch. Self-container and reality-question controls yielded zero viable pairs. This audit applies validity criteria developed by the same author; the framework is independently deposited so readers can verify ratings against the stated standards.

\paragraph{Verdict.}
We upgrade the claim from ``Causally Suggestive'' to ``Mechanistically Supported'' (Figure~\ref{fig:heatmap}). The subspaces replicate on a minor model version update (Figure~\ref{fig:replication}), beat shuffled-label and random-subspace baselines (measurement validity), are absent on non-belief tasks at L52/L53 (construct validity), generalize to false beliefs and observability-mediated knowledge (external validity), are framing-sensitive rather than framing-invariant (interpretive validity, Figure~\ref{fig:framing}), 19 of 20 positive effects are robust to unmeasured confounding (Figure~\ref{fig:evalues}), and the answer lookback dependency structure transfers cross-model to Qwen2.5-14B-Instruct (X2, Section~\ref{sec:crossmodel}). The remaining gap to ``Triangulated'' would require convergent evidence from an independent method (weight analysis, probing) and cross-model replication with proper multi-seed statistics.

\paragraph{Revised claim.}
We originally proposed contracting ``belief tracking via pointer dereference'' to ``attention-mediated retrieval.'' E12--E15 warrant a more specific characterization. The subspaces encode character-state associations that update with observability and state changes, accessible specifically through character-perspective queries. This is richer than ``retrieval'' but more precise than ``belief tracking''---belief tracking implemented as entity-state binding, not as a separate cognitive primitive. We propose the term \emph{belief-constitutive entity binding}: the entity bindings constitute beliefs in the functional sense, encoding what each character would report about each container.

\paragraph{Implications for mechanistic interpretability.}
Three of four dissociation experiments overturned their predictions, each time strengthening the original claim. Without pre-registration, these results could be dismissed as post-hoc rationalization. With it, the pattern of overturned predictions is itself evidence: we designed experiments to find entity binding and found belief tracking instead. E12, which we expected to undermine the paper's interpretation, provides the clearest evidence \emph{for} it.


\section*{Data and Code Availability}

All audit materials, pre-registration documents, and experiment code are available at \url{https://github.com/elliottower/lookback-validity-audit}.

\bibliographystyle{tmlr}
\bibliography{references}

\end{document}
